% Part: sets-functions-relations
% Chapter: functions
% Section: isomorphisms

\documentclass[../../../include/open-logic-section]{subfiles}

\begin{document}

\olfileid{sfr}{fun}{iso}
\olsection{آيزومورفيزم}

\begin{explain}
\emph{آيزومورفيزم} داسې دوه‌اړخيزه يو پر يو تابع ده چې د هغو سټونو جوړښت ساتي چې سره تړي۔ جوړښت د سټونو د !!{element}s تر منځ له اړيکو ټاکل کېږي۔ دا دوه سټونه $X=\{1,2,3\}$ او $Y=\{4,5,6\}$ په پام کښې ونيسئ۔ دواړه سټونه د ځاے‌ناستي، له ... نه کم او له ... نه زيات په اړيکو جوړښت مومي۔ د دواړو سټونو تر منځ آيزومورفيزم هغه !!a{bijection} ده چې همدا اړيکې ساتي۔ نو !!a{bijective} تابع $f \colon X
\to Y$ آيزومورفيزم ده که $i<j$ هغه او يوازې هغه وخت چې $f(i)<f(j)$، $i>j$ هغه او يوازې هغه وخت چې $f(i)>f(j)$، او $j$ د $i$ ځاے‌ناستے وي هغه او يوازې هغه وخت چې $f(j)$ د $f(i)$ ځاے‌ناستے وي۔
\end{explain}

\begin{defn}[آيزومورفيزم]
فرض کړئ $U$ جوړه $\langle X, R\rangle$ او $V$ جوړه $\langle
Y, S\rangle$ ده، چې $X$ او $Y$ سټونه دي او $R$ او $S$ په ترتيب پر $X$ او $Y$ اړيکې دي۔ له $X$ نه $Y$ ته !!{bijection} $f$ له $U$ نه $V$ ته \emph{آيزومورفيزم} ده هغه او يوازې هغه وخت چې اړيکيز جوړښت وساتي؛ يعنې په $X$ کښې د هر $x_{1}$ او $x_{2}$ دپاره $\tuple{x_1,x_2} \in R$ هغه او يوازې هغه وخت چې $\tuple{f(x_1),f(x_2)} \in S$۔
\end{defn}

\begin{ex}
دا دوه سټونه $X=\{1,2,3\}$ او $Y=\{4,5,6\}$ او د کمښت او زياتوالي اړيکې په پام کښې ونيسئ۔ تابع $f\colon X \to
Y$، چې $f(x) = 7-x$، د $\tuple{X,<}$ او $\tuple{Y,>}$ تر منځ آيزومورفيزم ده۔
\end{ex}

\end{document}
